% ============================================================== 5. Discussion
\section{Discussion}
\label{sec:discussion}

\subsection{Principal findings}

In Fed-ISIC2019, image-level splitting places sibling images of the same lesion in training and test sets
for 0.1--93\% of each client's test images in the official split. Models looked about 8 balanced-accuracy
points better on the full image-level test sets than on their new-lesion subsets with any of three encoder
configurations, and up to 16 points better for centralized training (RQ2). A randomized experiment on fixed test images confirmed that the sibling images themselves cause this gain (9.1 points for FedAvg) and that it grows with how closely a classifier fits its training lesions. Leakage left the
ranking of fine-tuned methods on the benchmark metric about as stable as re-running with another seed, but
it reordered methods on per-client accuracy and for a foundation-model encoder, where the most flexible
head moved from first to seventh place. The size of the apparent inflation was associated with how closely
a model can fit its training lesions rather than with personalization as such (RQ3). Calibration and
conformal coverage varied more between clients than between methods, and a single global conformal
threshold left some clients well below the target coverage (RQ1). A pre-specified fixed shrinkage rule
failed; an empirical-Bayes variant kept every client's mean coverage near the target and improved the
worst-client tail for one of two encoders (RQ4). Accuracy, calibration and transferred coverage were worse
on unseen clients and much worse on smartphone images.

\subsection{Why leakage inflates apparent accuracy}

Two observations constrain the explanation. First, frozen encoders with linear heads, which have far less
trainable capacity than a fully fine-tuned network, still showed about 8 points of apparent inflation.
Second, inflation was larger for classifiers that fit their training lesions more closely: linear heads
about 6--9 points, the two-layer FedPer-style head about 11, fully fine-tuned networks trained for 20 epochs
12--16, and a 1-nearest-neighbor rule about 22. One explanation consistent with both observations is that
images of one lesion lie close together in feature space, so a test image whose sibling is in training sits
next to a correctly labelled neighbor even in a generic ImageNet embedding, and that how much of this
proximity becomes apparent accuracy depends on how closely the classifier follows individual training
points. This account fits the large inflation of non-personalized pooled training, the small inflation of
FedBN and FedAvg-FT, and the reversal of the local-only advantage between model~A and model~B. The H3b
hypothesis, formulated in terms of personalization, was therefore supported for model~A, but the
supporting difference came mainly from local-only training. Because the capacity comparisons span
different optimizers, budgets and head architectures, and the nearest-neighbor analysis is exploratory,
we regard this as an interpretation, not as a tested causal mechanism. The two additional experiments strengthen this account for the
comparisons they control. In E8, the exposure effect on the same randomized target images ordered the
methods in the same way as the observational contrast (local-only $>$ pooled $>$ FedPer $>$ FedAvg for
model~A; the two-layer head far above the linear heads for model~B), and in E9 a FedAvg model with the
FedPer-style two-layer head inflated as much as the FedPer-style head itself. Head capacity, not the
client-specific layer, therefore explains the behaviour of the FedPer-style head. For the fine-tuned
network, the methods still differ in optimizer and training budget, so fitting capacity remains the most
consistent explanation rather than an isolated cause.

Freezing a foundation model did not remove the apparent inflation: PanDerm showed as much as an ImageNet
encoder under FedAvg, and its most flexible head was the variant most favored by the leaky split. Since
PanDerm's pretraining data are not fully public, some of its specific behavior could reflect overlap with
the evaluation images. The A-frozen control shows that inflation of this size does not require such
overlap, but it does not tell how much of PanDerm's own inflation is due to it.

\subsection{Implications for federated benchmarking}

Differences between FL methods on Fed-ISIC2019 are often a few balanced-accuracy points: in FLamby's
benchmark and in the evaluation of Tavakoli et al., FedAvg, FedProx, SCAFFOLD and personalized methods lie
within about 5 points of each other \cite{flamby2022,tavakoli2024fendafl}. The apparent leakage effect is
of the same size or larger and is unevenly distributed, with most of it in one client. Two consequences
follow. Absolute accuracies reported on image-level splits overstate performance on new lesions, and
leakage can change comparative performance and rankings, with effects that depend on the methods and the
metric. We observed both a benign and a harmful case: for fine-tuned EfficientNet-B0 the pooled ranking was
stable, whereas for frozen PanDerm features it was not. Because the harmful case cannot be recognized
without a grouped split, we recommend lesion- or patient-grouped splits as the default for FL benchmarks
with repeated imaging. This extends previous findings on duplicates and problematic partitions in
centralized dermatology datasets \cite{cassidy2022isic,groger2023reliable,abhishek2025dermamnist} and on
subject-level leakage in other imaging domains \cite{bussola2021tiles,yagis2021leakage} to a federated
benchmark with acquisition-source clients, and to calibration and coverage. The lesion-grouped splits used
here are released with the code.

\subsection{Reliability at the level of the site}

For a clinician, the relevant question is whether probabilities and prediction sets are reliable at their
own site. Four results bear on this. First, pooled ECE values of 0.02--0.05 concealed per-client values up
to 0.26, so pooled calibration metrics are not informative about individual sites, although small clients
also have imprecise ECE. Second, the apparent difference in ECE between new-lesion and full leaky test sets was explained by
sample size and composition; this illustrates that binned ECE should be compared only at matched sample
sizes, and that NLL and Brier score are useful companions. Third, temperature scaling fitted on lesion-grouped validation
lesions made the most strongly fitted models worse on a leaky test set, so calibration studies that tune on
one kind of split and test on another can reach opposite conclusions; validation and test data should be
grouped in the same way. Fourth, a global threshold, including FCP's \cite{lu2023fcp}, targets coverage on
the mixture of clients and under-covered clients whose scores were shifted, whereas client-wise thresholds
were correct on average but noisy in small clients. Client-wise thresholds rest on exchangeability within a
client, which repeated images of one lesion weaken; randomizing whole lesions between calibration and
evaluation and the one-image-per-lesion analysis address this empirically, while hierarchical conformal
methods offer formal guarantees for grouped data \cite{dunn2022hierarchical,lee2023hierarchical}.

EB-ShrinkFCP addresses the noise of client-wise thresholds without returning to a single shared threshold.
Its weight $w_c=\tau^2/(\tau^2+\sigma_c^2)$ estimates from the data how much the clients differ, so it
borrows strength only when they agree; the fixed weight of ShrinkFCP borrowed regardless and imported the
global threshold's bias into the shifted clients. The attainable gain is limited by evaluation noise: with
about 44 evaluation images per draw in the smallest client, even a rule with exactly nominal conditional
coverage has a 10th percentile of observed coverage well below 0.90, although trivial rules such as
returning all classes would reach full coverage at the cost of useless sets. The evidence on RQ4 is
therefore mixed. Across the tested models and clients, EB-ShrinkFCP kept mean client coverage close to the
target; it met its pre-specified criterion on held-out models for model~A, with an interval that excludes
zero but also includes gains below 2 points; no clear improvement was detected for model~B, so the joint criterion
was not met; and class-conditional coverage was not established. The post hoc analyses point the same way but do not strengthen the claim: the gain is visible in the
variability of thresholds, but with one image per lesion its interval includes zero, and its size is
comparable to the Monte Carlo variability of a 100-draw q10. The rule exchanges two
aggregate statistics per client, a quantile and a variance, rather than images or score lists; we did not
analyze its formal privacy properties. Related work on federated conformal prediction has addressed label
shift \cite{plassier2023labelshift}, one-shot aggregation \cite{humbert2023oneshot}, personalization
through localization \cite{min2025pfcp} and group-conditional coverage \cite{wen2026gcfcp}; of these we
compared only a label-shift-weighted threshold.

Under the smartphone-image shift, transferred thresholds lost 23--30 points of coverage, consistent with the
absence of any guarantee under distribution shift \cite{tibshirani2019covariate,barber2023beyond}. Local
recalibration restored marginal coverage only by returning four to five of six classes, and coverage for
actinic keratosis and squamous cell carcinoma stayed below the target. Recalibration can make a model's
uncertainty honest, but it cannot make an unsuitable model useful, and marginal coverage can hide
class-conditional failures \cite{lu2022fair}.

\subsection{Recommendations}

Based on these results, we suggest that federated dermatology studies (i) split data by lesion or patient
and report the leakage rate of any image-level split they compare with; (ii) report balanced accuracy,
calibration and coverage per client, not only pooled, together with the uncertainty of small-client
estimates; (iii) fit temperature and conformal thresholds on data grouped in the same way as the test data;
(iv) state which coverage target a conformal rule serves (mixture, client or class) and compare locally
calibrated and pooled thresholds for that target before adopting a shrinkage rule; (v) evaluate on held-out
clients and, before deployment in a new imaging modality, recalibrate locally and report set sizes and
per-class coverage; and (vi) fix primary analyses in advance and report negative results, in line with
general reporting guidance for clinical prediction models \cite{collins2024tripodai}.

\subsection{Value of a fixed protocol and negative results}

Of the seven pre-specified tests, four were met in full (Gates 1 and 2, the PAD-UFES-20 transfer rule and E8), two were met for one of two encoders (H3b and 5b), and one failed (RQ4). Without a protocol, the failure of ShrinkFCP could have been
followed by adjusting its constant until it succeeded on the same models. Instead, the failure motivated a
new rule that was specified before the held-out models existed and tested on them, and the post hoc results
are labelled as such. The protocol also prevented a selective reading of H3b: the hypothesis was supported
for model~A but not for model~B. The protocol was version-controlled rather than registered externally, and
two amendments were written after results were known (Supplementary~S9); we therefore present the chronology in
full rather than describing the study as pre-registered in the strict sense
\cite{nosek2018preregistration,kapoor2023leakage,varoquaux2022failures}.

\subsection{Limitations}

This study has several limitations. It audits one benchmark whose six clients are acquisition sources,
three of them from one institution; other federated datasets with patient or lesion identifiers are needed
to test generality. Patient identifiers are not available, so different lesions of one patient can fall on
both sides of a lesion-grouped split, and images of one lesion remain correlated within the calibration and
evaluation halves. Clients 4 and 5 have fewer than 200 test images, so their per-client estimates are
imprecise, and much of the leakage effect is driven by client~0. The apparent inflation compares different
groups of test images and can reflect differences in difficulty and multiplicity even after
standardization; the randomized experiment E8 addresses this for lesions with several images, but not for single-image lesions, and its method comparisons still differ in optimizer and budget. Binned ECE is biased in small samples, which affected our original leaky-versus-new
contrast. Three to five seeds limit the precision of ranking comparisons, and the bootstrap intervals
reflect test-sample uncertainty given the trained models, not retraining variability. The pooled and
local-only references use a different optimizer and budget from the federated methods, FedAvg-FT used a
single pre-specified number of fine-tuning steps, and the benchmark clients and test pools were reused
across method development and evaluation. Overlap between PanDerm's pretraining data and the evaluation
images could not be ruled out. Our FedAvg reproduction was 4.2 points below one published value, though
above FLamby's own results. EB-ShrinkFCP met its criterion for one of two encoders, it does not carry a
finite-sample per-client guarantee, and its variance approximation for clustered images is heuristic.
Finally, the mechanism analysis is exploratory.

\subsection{Future work}

Future work should extend the analysis to multi-site datasets with patient identifiers, examine
hierarchical conformal methods that target new patients or lesions directly, compare EB-ShrinkFCP with
personalized, label-shift-aware and group-conditional federated conformal methods
\cite{min2025pfcp,plassier2023labelshift,wen2026gcfcp} under the same coverage target, and target
class-conditional coverage for rare and commonly confused classes.
